\section{Removal Game}
\label{sec:cses-removal-game}

\problemstatement{An array of $n$ numbers is on a table. Two players
alternate; on each turn a player removes either the leftmost or rightmost
element and adds its value to their score. Both play optimally to maximise
their own score. Find the first player's final score.}

\begin{patternbox}
\emph{``Two players take from the ends, both optimal''} is an interval
game DP. Let $\textit{dp}[l][r]$ be the best score \emph{difference} (current
player minus opponent) achievable on the subarray $[l,r]$; the current
player takes an end and then \emph{becomes} the opponent on the rest.
\end{patternbox}

\subsection*{State and transition via score difference}

Track $\textit{dp}[l][r]$ = the maximum of (current player's score $-$
opponent's score) over the subarray $a[l..r]$, assuming optimal play. The
current player takes $a[l]$ or $a[r]$; after that the opponent plays
optimally on the remaining interval, so their advantage is subtracted:
\[
  \textit{dp}[l][r]=\max\!\big(a[l]-\textit{dp}[l+1][r],\
  a[r]-\textit{dp}[l][r-1]\big),\qquad \textit{dp}[l][l]=a[l].
\]
If $\textit{total}=\sum a_i$ and $d=\textit{dp}[0][n-1]$ is the final
difference, the first player's score is $(\textit{total}+d)/2$.

\begin{ideabox}[Score Difference Turns Two Objectives into One]
Modelling the game by the \emph{difference} folds ``maximise mine'' and
``opponent minimises mine'' into a single recurrence: whatever the
opponent gains counts against me, so I subtract their optimal advantage on
the remainder. This zero-sum trick avoids storing two separate scores.
\end{ideabox}

\subsection*{Algorithm}

\begin{enumerate}
  \item $\textit{dp}[i][i]=a[i]$.
  \item For increasing interval length, $\textit{dp}[l][r]=\max(a[l]-
        \textit{dp}[l+1][r],\ a[r]-\textit{dp}[l][r-1])$.
  \item First player's score $=(\textit{total}+\textit{dp}[0][n-1])/2$.
\end{enumerate}

\subsection*{Dry run}

\dryrun{$a=[4,5,1,3]$, total $=13$.}

\begin{itemize}
  \item Length 1: $\textit{dp}[i][i]=a[i]$.
  \item Length 2: $\textit{dp}[0][1]=\max(4-5,5-4)=1$;
        $\textit{dp}[2][3]=\max(1-3,3-1)=2$; $\textit{dp}[1][2]=\max(5-1,
        1-5)=4$.
  \item Full: $\textit{dp}[0][3]=\max(4-\textit{dp}[1][3],
        3-\textit{dp}[0][2])$. With $\textit{dp}[1][3]=\max(5-2,3-4)=3$
        and $\textit{dp}[0][2]=\max(4-4,1-1)=0$: $\max(4-3,3-0)=3$.
  \item First player's score $=(13+3)/2=8$.
\end{itemize}

\subsection*{C++ implementation}

\begin{lstlisting}
#include <bits/stdc++.h>
using namespace std;

int main() {
    int n;
    cin >> n;
    vector<long long> a(n);
    long long total = 0;
    for (auto& x : a) { cin >> x; total += x; }

    vector<vector<long long>> dp(n, vector<long long>(n, 0));
    for (int i = 0; i < n; i++) dp[i][i] = a[i];

    for (int len = 2; len <= n; len++)
        for (int l = 0; l + len - 1 < n; l++) {
            int r = l + len - 1;
            dp[l][r] = max(a[l] - dp[l + 1][r], a[r] - dp[l][r - 1]);
        }

    cout << (total + dp[0][n - 1]) / 2 << "\n";
    return 0;
}
\end{lstlisting}

\begin{complexitybox}
\textbf{Time Complexity.} $O(n^2)$ -- one $O(1)$ transition per interval.
\textbf{Space Complexity.} $O(n^2)$ for the interval table.
\end{complexitybox}

\begin{takeawaybox}
For symmetric two-player, zero-sum games, track the \emph{score
difference} from the mover's perspective and subtract the opponent's
optimal remainder. This collapses adversarial optimisation into a single
maximisation -- the standard move for end-taking and Nim-like game DPs.
\end{takeawaybox}
